\documentclass[10pt,a4paper]{amsart}
\usepackage{amsmath,amssymb,amsthm,graphicx}
\usepackage{verbatim}
\newenvironment{reference}{\comment}{\endcomment}
%\newenvironment{reference}{}{}
\newenvironment{slogan}{\comment}{\endcomment}
\newenvironment{history}{\comment}{\endcomment}

% For commutative diagrams we use Xy-pic
\usepackage[all]{xy}

% We use 2cell for 2-commutative diagrams.
\xyoption{2cell}
% Theorem environments.
%
\theoremstyle{plain}
\newtheorem{theorem}[subsection]{Theorem}
\newtheorem{proposition}[subsection]{Proposition}
\newtheorem{lemma}[subsection]{Lemma}

\theoremstyle{definition}
\newtheorem{definition}[subsection]{Definition}
\newtheorem{example}[subsection]{Example}
\newtheorem{exercise}[subsection]{Exercise}
\newtheorem{situation}[subsection]{Situation}

\theoremstyle{remark}
\newtheorem{remark}[subsection]{Remark}
\newtheorem{remarks}[subsection]{Remarks}

\numberwithin{equation}{subsection}

% Macros
%
\def\lim{\mathop{\mathrm{lim}}\nolimits}
\def\colim{\mathop{\mathrm{colim}}\nolimits}
\def\Spec{\mathop{\mathrm{Spec}}}
\def\Hom{\mathop{\mathrm{Hom}}\nolimits}
\def\Ext{\mathop{\mathrm{Ext}}\nolimits}
\def\SheafHom{\mathop{\mathcal{H}\!\mathit{om}}\nolimits}
\def\SheafExt{\mathop{\mathcal{E}\!\mathit{xt}}\nolimits}
\def\Sch{\mathit{Sch}}
\def\Mor{\mathop{\mathrm{Mor}}\nolimits}
\def\Ob{\mathop{\mathrm{Ob}}\nolimits}
\def\Sh{\mathop{\mathit{Sh}}\nolimits}
\def\NL{\mathop{N\!L}\nolimits}
\def\CH{\mathop{\mathrm{CH}}\nolimits}
\def\proetale{{pro\text{-}\acute{e}tale}}
\def\etale{{\acute{e}tale}}
\def\QCoh{\mathit{QCoh}}
\def\Ker{\mathop{\mathrm{Ker}}}
\def\Im{\mathop{\mathrm{Im}}}
\def\Coker{\mathop{\mathrm{Coker}}}
\def\Coim{\mathop{\mathrm{Coim}}}

% Boxtimes
%
\DeclareMathSymbol{\boxtimes}{\mathbin}{AMSa}{"02}

%
% Macros for moduli stacks/spaces
%
\def\QCohstack{\mathcal{QC}\!\mathit{oh}}
\def\Cohstack{\mathcal{C}\!\mathit{oh}}
\def\Spacesstack{\mathcal{S}\!\mathit{paces}}
\def\Quotfunctor{\mathrm{Quot}}
\def\Hilbfunctor{\mathrm{Hilb}}
\def\Curvesstack{\mathcal{C}\!\mathit{urves}}
\def\Polarizedstack{\mathcal{P}\!\mathit{olarized}}
\def\Complexesstack{\mathcal{C}\!\mathit{omplexes}}
% \Pic is the operator that assigns to X its picard group, usage \Pic(X)
% \Picardstack_{X/B} denotes the Picard stack of X over B
% \Picardfunctor_{X/B} denotes the Picard functor of X over B
\def\Pic{\mathop{\mathrm{Pic}}\nolimits}
\def\Picardstack{\mathcal{P}\!\mathit{ic}}
\def\Picardfunctor{\mathrm{Pic}}
\def\Deformationcategory{\mathcal{D}\!\mathit{ef}}
\usepackage{cleveref}
\crefname{originalreference}{original reference}{original references}
\Crefname{originalreference}{Original reference}{Original references}
\makeatletter
\newlabel{lemma-map-completions-etale}{{0GI0}{source}{}{sourcepointer-1}{}}
\newlabel{lemma-map-completions-etale@cref}{{[originalreference][0][]0GI0}{[1][0][]source}{}{}{}}
\makeatother
\crefname{originallemma}{original lemma}{original lemmas}
\Crefname{originallemma}{Original lemma}{Original lemmas}
\makeatletter
\newlabel{spaces-more-morphisms-lemma-check-smoothness-on-infinitesimal-nbhds}{{0APQ}{source}{}{sourcepointer-2}{}}
\newlabel{spaces-more-morphisms-lemma-check-smoothness-on-infinitesimal-nbhds@cref}{{[originallemma][0][]0APQ}{[1][0][]source}{}{}{}}
\makeatother
\crefname{originaldefinition}{original definition}{original definitions}
\Crefname{originaldefinition}{Original definition}{Original definitions}
\makeatletter
\newlabel{spaces-properties-definition-reduced-induced-space}{{047X}{source}{}{sourcepointer-3}{}}
\newlabel{spaces-properties-definition-reduced-induced-space@cref}{{[originaldefinition][0][]047X}{[1][0][]source}{}{}{}}
\makeatother
\crefname{originalreference}{original reference}{original references}
\Crefname{originalreference}{Original reference}{Original references}
\makeatletter
\newlabel{theorem-dilatations-general}{{0ARB}{source}{}{sourcepointer-4}{}}
\newlabel{theorem-dilatations-general@cref}{{[originalreference][0][]0ARB}{[1][0][]source}{}{}{}}
\makeatother
\title{Etaleness transfers between a dilatation and its formal model}
\author{The Stacks Project Authors}
\date{stacks-project:restricted.tex:local-snapshot-20261012; GNU FDL 1.2 or later}
\begin{document}
\maketitle
\noindent\textit{Editorial note.} All delivered mathematical text is the author's own, copied from the pinned Stacks Project chapter file restricted.tex; no statement or proof step was rewritten, reordered or removed.\par\medskip
\noindent\textit{Editorial note.} The counted proof is the author's complete argument: the proof handles the open part by the isomorphism over U, writes X\_\{/T\} as the colimit of the infinitesimal neighbourhoods, deduces etaleness of each base change from g, and applies the infinitesimal etaleness criterion\par\medskip
\noindent\textit{Editorial note.} The label lemma-output-etale names the counted statement, and the proof environment immediately following it in the source is delivered verbatim as its proof.\par\medskip
\noindent\textit{Editorial note.} The counted statement is carried by the author's own lemma environment in the chapter source.\par\medskip
\noindent\textit{Allowed standard prerequisites (editorial).}\begin{itemize}
\item the dilatations theorem (theorem-dilatations-general of the same chapter) and the correspondence between f : X' -> X and g : W -> X\_\{/T\}
\item etale f gives etale formal completion (lemma-map-completions-etale of the same chapter)
\item the reduced induced closed subspace Z on T, its infinitesimal neighbourhoods Z\_n, and X\_\{/T\} = colim Z\_n
\item the identification of f\textasciicircum{}\{-1\}(Z\_n) with X' x\_X Z\_n and etaleness of its map to Z\_n from the assumed etaleness of g
\item etaleness checked on infinitesimal neighbourhoods (More on Morphisms of Spaces, Lemma spaces-more-morphisms-lemma-check-smoothness-on-infinitesimal-nbhds)
\end{itemize}\medskip
\ifdefined\hypertarget\hypertarget{sourcepointer-1}{}\fi
\noindent\textit{Source reference.} Original reference 0GI0: sources/stacks-restricted/pinned-original-source.tex, source lines 6026--6026. Prerequisite pointer invoked at proof line(s) 7345; it names a standard fact of the same chapter that the counted argument uses, and every step of the counted proof built on it is retained verbatim in the delivered body.\par\medskip
\ifdefined\hypertarget\hypertarget{sourcepointer-2}{}\fi
\noindent\textit{Source reference.} Original lemma 0APQ: sources/stacks-chow/tags.txt, source lines 13483--13483. Prerequisite pointer invoked at proof line(s) 7359; it names a standard fact of a sibling chapter that the counted argument uses, and every step of the counted proof built on it is retained verbatim in the delivered body.\par\medskip
\ifdefined\hypertarget\hypertarget{sourcepointer-3}{}\fi
\noindent\textit{Source reference.} Original definition 047X: sources/stacks-chow/tags.txt, source lines 5524--5524. Prerequisite pointer invoked at proof line(s) 7353; it names a standard fact of a sibling chapter that the counted argument uses, and every step of the counted proof built on it is retained verbatim in the delivered body.\par\medskip
\ifdefined\hypertarget\hypertarget{sourcepointer-4}{}\fi
\noindent\textit{Source reference.} Original reference 0ARB: sources/stacks-restricted/pinned-original-source.tex, source lines 7015--7015. Prerequisite pointer invoked at proof line(s) 7338; it names a standard fact of the same chapter that the counted argument uses, and every step of the counted proof built on it is retained verbatim in the delivered body.\par\medskip
\begin{lemma}
\label{lemma-output-etale}
With assumptions and notation as in Theorem \ref{theorem-dilatations-general}
let $f : X' \to X$ correspond to $g : W \to X_{/T}$.
Then $f$ is \'etale if and only if $g$ is \'etale.
\end{lemma}

\begin{proof}
If $f$ is \'etale, then $g$ is \'etale by
Lemma \ref{lemma-map-completions-etale}.
Conversely, assume $g$ is \'etale.
Since $f$ is an isomorphism over $U$ we see that $f$ is \'etale
over $U$. Thus it suffices to prove that $f$ is \'etale
at any point of $X'$ lying over $T$. Denote $Z \subset X$
the reduced closed subspace whose underlying topological
space is $|Z| = T \subset |X|$, see
Properties of Spaces, Definition
\ref{spaces-properties-definition-reduced-induced-space}.
Letting $Z_n \subset X$ be the $n$th infinitesimal neighbourhood
we have $X_{/T} = \colim Z_n$. Since $X'_{/T} = W \to X_{/T}$
we conclude that $f^{-1}(Z_n) = X' \times_X Z_n \to Z_n$
is \'etale by the assumed \'etaleness of $g$.
By More on Morphisms of Spaces, Lemma
\ref{spaces-more-morphisms-lemma-check-smoothness-on-infinitesimal-nbhds}
we conclude that $f$ is \'etale at points lying over $T$.
\end{proof}

\end{document}
